% ============================================================
%  INTRODUCCIÓN — Tesis de desarrollo
%  Tema: Reconocimiento de voz, traducción automática y síntesis
%        de voz aplicados al quechua para la comunicación
%        médico-paciente.
%  Estructura: Sistema de 6 bloques (Seminario de Investigación II,
%        Msc. Raul Huillca Huallparimachi).
%
%  Requiere (biblatex-apa / APA7):
%     \usepackage[style=apa,backend=biber]{biblatex}
%     \addbibresource{referencias.bib}
%     \addbibresource{41articulos_Scopus.bib}
%
%  ⚠ ANTES DE COMPILAR: en 41articulos_Scopus.bib corrige la clave
%     malformada  @ARTICLE{Huamán Curo2026167,  -> {HuamanCuro2026,}
%     (tiene un espacio y una tilde; BibTeX no admite espacios en la
%     clave). Aquí se cita como \textcite{HuamanCuro2026}.
%  ⚠ Hay entradas duplicadas entre ambos .bib (p. ej. Ortega 2020:
%     ortega2020quechua == Ortega2020325). Se usa la curada.
% ============================================================

\section*{Introducción}

% ---------- BLOQUE 1: Contexto y problema general ----------
El procesamiento del lenguaje natural y, en particular, las tecnologías del
habla —reconocimiento automático de voz (ASR), traducción automática (MT) y
síntesis de voz (TTS)— constituyen un área de la inteligencia artificial
orientada a que las personas puedan interactuar con los sistemas informáticos
en su lengua materna. Los avances recientes en aprendizaje profundo han
permitido construir modelos multilingües de gran escala, como Whisper
\parencite{radford2023whisper}, MMS \parencite{pratap2024mms} y NLLB
\parencite{costajussa2022nllb}, que alcanzan niveles elevados de desempeño en
reconocimiento, traducción y generación de voz. Sin embargo, estas capacidades
se concentran en un reducido conjunto de lenguas de alto recurso, para las que
existen grandes volúmenes de datos transcritos y corpus paralelos; las lenguas
de bajo recurso quedan sistemáticamente al margen, situación que se agrava en
lenguas de morfología aglutinante y polisintética como el quechua, cuya
segmentación morfológica dificulta el entrenamiento de los modelos
\parencite{ortega2020quechua}. En el Perú, el quechua es hablado por alrededor
de 3,8 millones de personas \parencite{inei2017}, muchas de ellas monolingües en
zonas rurales andinas; esta brecha tecnológica se traduce en barreras
lingüísticas concretas en la atención de salud, donde los pacientes
quechuahablantes enfrentan dificultades para comunicarse con el personal
asistencial que no domina la lengua \parencite{begazo2026,cuaila2022}. Estas
barreras comprometen la calidad del diagnóstico, el consentimiento informado y
la adherencia al tratamiento. Resulta
pertinente, por tanto, investigar soluciones tecnológicas del habla que apoyen
la comunicación médico-paciente en quechua.

% ---------- BLOQUE 2: Tecnología o método ----------
Para abordar el reconocimiento de voz en quechua se han empleado desde modelos
ocultos de Markov \parencite{Zevallos202061} hasta enfoques \textit{end-to-end}
y de aprendizaje autosupervisado basados en representaciones preentrenadas
\parencite{Chen20246444}. En traducción automática, la evolución va de los
métodos estadísticos a la traducción neuronal con arquitecturas Transformer y el
ajuste fino de modelos multilingües preentrenados
\parencite{Ahmed2023156,Barbaran2026464,Dhawan2026119}. En síntesis de voz, se
han evaluado arquitecturas como Tacotron~2, FastPitch, VITS y XTTS para el
quechua chanca \parencite{HuamanCuro2026}. La presente investigación se
fundamenta en la integración de estos tres componentes en una arquitectura en
cascada (ASR${\rightarrow}$MT${\rightarrow}$TTS), de modo que un enunciado
hablado en quechua pueda reconocerse, traducirse al español y, en sentido
inverso, sintetizarse nuevamente como voz en quechua, habilitando un canal de
diálogo bidireccional; este tipo de encadenamiento ya ha sido explorado para el
quechua, aunque con fines de subtitulado
\parencite{Alvarez-Crespo2023837}; sin embargo, ninguno de estos encadenamientos
se ha orientado ni validado en el dominio de la salud.

La evidencia previa muestra resultados
prometedores por componente: un ASR con HTK alcanzó 12,70\,\% de WER
\parencite{Zevallos202061} y las técnicas de aumento de datos lo mejoraron hasta
22,75\,\% de WER en condiciones de bajo recurso \parencite{Zevallos20223518};
los sistemas de traducción de voz QUESPA alcanzaron entre 16 y 26,7 BLEU en las
campañas IWSLT \parencite{Ortega2023261,Ortega2024173,Ortega2025260}; y los
sintetizadores TTS reportaron valores de MOS de hasta 4,9
\parencite{HuamanCuro2026}. Recursos como el corpus Huqariq
\parencite{zevallos2022huqariq} y el modelo QuBERT \parencite{Zevallos20221} han
contribuido a mitigar la escasez de datos. No obstante, persisten limitaciones
relevantes: los datos disponibles son escasos y heterogéneos entre variantes
dialectales (Cusco-Collao, Ayacucho-Chanca), no existen corpus orientados al
dominio médico, los componentes se evalúan de forma aislada y las pruebas se
realizan en dominios de propósito general, sin validación en escenarios clínicos
reales.

% ---------- BLOQUE 3: Vacío de investigación ----------
La literatura revisada evidencia que el reconocimiento, la traducción y la
síntesis de voz para el quechua han sido abordados con resultados alentadores,
principalmente en el marco de competencias como AmericasNLP e IWSLT
\parencite{Oncevay2021194,Ebrahimi2024236}. Sin embargo, la mayoría de los
trabajos: (i) optimiza y evalúa un único componente del flujo; (ii) opera en
dominios de propósito general —textos literarios, noticias o actas—; y (iii) se
concentra en una sola variante dialectal. En consecuencia, no se dispone de
evidencia experimental sobre un sistema integrado ASR+MT+TTS para el quechua
diseñado y validado específicamente para la comunicación médico-paciente, ni de
corpus de dominio sanitario, ni de mediciones de utilidad en condiciones de uso
clínico. Esta investigación delimita su vacío en la integración de los tres
componentes y su evaluación conjunta —desempeño y usabilidad— en el contexto
médico-paciente, para una variante dialectal definida (quechua sureño/Collao de
la región Cusco).

% ---------- BLOQUE 4: Contexto y problema específico ----------
El problema específico se relaciona con la necesidad de establecer el
comportamiento de una arquitectura integrada en cascada que combine ASR, MT y
TTS para habilitar el diálogo quechua-español en entornos de salud. En
particular, es necesario determinar si su implementación permite alcanzar un
desempeño aceptable en cada componente —medido con WER para el reconocimiento,
BLEU y chrF para la traducción \parencite{papineni2002bleu,popovic2015chrf}, y
MOS para la síntesis— y cómo se compara con las líneas base existentes, así como
caracterizar los recursos y la latencia que introduce su ejecución en un
escenario de consulta médica. La evaluación requiere, además, delimitar el
corpus de frases del dominio sanitario, la variante dialectal y las condiciones
del escenario de uso.

\paragraph{Pregunta de investigación.}
¿Cómo se comporta experimentalmente una arquitectura integrada en cascada de
reconocimiento de voz, traducción automática y síntesis de voz para el par
quechua-español, en términos de desempeño por componente (WER, BLEU/chrF, MOS),
latencia y usabilidad, aplicada a la comunicación médico-paciente, frente a
líneas base definidas bajo condiciones comparables?

% ---------- BLOQUE 5: Propuesta, objetivo y justificación ----------
En esta investigación se propone diseñar, desarrollar y evaluar un sistema
integrado de tecnologías del habla que articule ASR, MT y TTS para el par
quechua-español, orientado a apoyar la comunicación entre el personal de salud y
los pacientes quechuahablantes. El sistema se construirá mediante el ajuste fino
de modelos preentrenados de habla y traducción sobre datos de quechua,
complementados con un corpus de frases del dominio médico, y se evaluará
experimentalmente mediante comparaciones con métodos de referencia y métricas
por componente.

\paragraph{Objetivo general.}
Desarrollar y evaluar experimentalmente un sistema integrado de reconocimiento
de voz, traducción automática y síntesis de voz para el par quechua-español,
basado en el ajuste fino de modelos de aprendizaje profundo, con la finalidad de
apoyar la comunicación médico-paciente, considerando el desempeño por componente
(WER, BLEU/chrF, MOS), la latencia y la usabilidad frente a líneas base
definidas.

\paragraph{Objetivos específicos.}
\begin{itemize}
  \item Construir un corpus quechua-español de dominio médico para el
        entrenamiento y la evaluación del sistema.
  \item Implementar y ajustar el módulo de reconocimiento de voz (ASR) para el
        quechua.
  \item Implementar y ajustar el módulo de traducción automática (MT)
        quechua-español bidireccional.
  \item Implementar y ajustar el módulo de síntesis de voz (TTS) para el quechua.
  \item Integrar los componentes en una arquitectura en cascada y evaluar su
        desempeño, latencia y usabilidad en un escenario de comunicación
        médico-paciente.
\end{itemize}

\paragraph{Justificación.}
Científicamente, la investigación aporta evidencia sobre la integración de
tecnologías del habla en una lengua de bajo recurso y morfología compleja.
Metodológicamente, establece un procedimiento de evaluación conjunta por
componente y de usabilidad en un dominio aplicado, además de un corpus
reutilizable. Tecnológicamente, permite determinar las capacidades, las
limitaciones y los costos de un sistema de este tipo. Socialmente, contribuye a
reducir las barreras lingüísticas en la atención de salud y a la revitalización
del quechua, en línea con el derecho de los pacientes a recibir atención en su
lengua materna \parencite{essalud2023}.

\paragraph{Metodología (anticipación).}
La investigación adoptará un enfoque cuantitativo y experimental: se emplearán
conjuntos de datos de quechua y un corpus de dominio médico, se ajustarán
modelos preentrenados y se compararán mediante métricas definidas. Los detalles
de los parámetros, los equipos y los procedimientos se desarrollarán en el
capítulo metodológico.

% ---------- BLOQUE 6: Organización de la tesis ----------
\paragraph{Organización de la tesis.}
La tesis se organiza en cinco capítulos. El \textbf{Capítulo~1} presenta el
planteamiento del problema, las preguntas de investigación, los objetivos y la
justificación. El \textbf{Capítulo~2} desarrolla los antecedentes y el marco
teórico, a partir de una revisión sistemática de literatura sobre ASR, MT y TTS
para el quechua y las lenguas de bajo recurso. El \textbf{Capítulo~3} describe la
metodología, la arquitectura propuesta, el corpus, el entorno experimental y los
procedimientos de evaluación. El \textbf{Capítulo~4} presenta los resultados
obtenidos y su análisis comparativo. El \textbf{Capítulo~5} expone las
conclusiones, las limitaciones de la investigación y las recomendaciones para
trabajos futuros.
